\section{Background} 

\iffalse
\subsection{Cyber-Physical Systems}
A \gls{cps} combines computational components with physical processes. Sensors provide information about the physical environment to the computational system, which processes the measurements and generates control actions that affect the physical system. \glspl{uav} are a representative example of \gls{cps} because their behavior emerges from the continuous interaction between physical sensors, state estimation, control algorithms, communication systems, and the aircraft itself. In this context, the integrity of sensor measurements is particularly important. A manipulation affecting a sensor input can propagate through the state-estimation and control chain and ultimately influence the physical trajectory of the \gls{uav}.
\fi

\subsection{Autonomous UAV Navigation}
An autonomous \gls{uav} must continuously estimate its position, velocity, attitude and other state variables in order to follow a predefined trajectory. During autonomous operation, the flight controller compares the estimated state of the vehicle with the desired state generated by the mission or path-following algorithm. The controller then generates appropriate control commands to reduce the difference between the current and desired states. For waypoint-based navigation, the desired trajectory is defined by a sequence of geographical coordinates. The \gls{uav} attempts to reach each waypoint in sequence before continuing with the following mission item. ArduCopter provides an autonomous \texttt{AUTO} mode in which the vehicle can execute missions containing operations such as take-off, waypoint navigation, and landing. The original Tractor Beam study also considers ArduCopter's autonomous waypoint navigation and describes how the vehicle uses its estimated position to follow the mission trajectory \cite{noh2019tractorbeam}.

\subsection{Flight Controller} The flight controller is the computational component responsible for estimating the UAV state and generating the control commands required to maintain the desired flight behavior. In this work, the flight controller is provided by ArduPilot ArduCopter running in Software-In-The-Loop (SITL). The use of SITL allows the complete flight-control stack to be executed without requiring physical UAV hardware. The simulated environment also provides access to internal telemetry and simulation variables, making it possible to compare the physical simulated state with the measurements available to the flight controller. 

\subsection{GPS and GNSS} Global Navigation Satellite Systems provide positioning and timing information to receivers located on or near the Earth's surface. GPS is one of the GNSS systems and is widely used by autonomous UAVs as a source of absolute position information. A GPS receiver estimates the position of the vehicle from measurements obtained from multiple satellites. The resulting latitude, longitude, altitude, velocity, and accuracy information is then provided to the navigation system. The use of GPS as an external source of absolute position is particularly important for autonomous flight because inertial sensors alone accumulate errors over time. The Tractor Beam paper describes GPS as a satellite-based navigation system in which receivers estimate their position, velocity, and time from satellite information and pseudorange measurements \cite{noh2019tractorbeam}. 

\subsection{Inertial Measurement Unit} An Inertial Measurement Unit (IMU) provides measurements from sensors such as accelerometers and gyroscopes. Unlike GPS, inertial measurements do not directly provide an absolute geographical position. Instead, the flight controller can integrate accelerations and angular rates to estimate changes in the vehicle state. However, inertial measurements contain noise and systematic errors. Consequently, the position estimate obtained exclusively from the IMU drifts over time. For this reason, autonomous UAVs commonly combine GPS and IMU measurements using sensor-fusion algorithms. 

\subsection{Extended Kalman Filter} ArduCopter uses an Extended Kalman Filter (EKF) to estimate the state of the vehicle by combining measurements from different sensors. The EKF continuously predicts the vehicle state using inertial measurements and compares the prediction with external measurements such as GPS. The difference between the predicted state and the GPS measurement is known as the innovation. If the innovation becomes inconsistent with the expected uncertainty of the system, the measurement may be rejected. The Tractor Beam study specifically analyzes the ArduCopter EKF and describes how position and velocity innovations are used to detect inconsistencies between GPS and inertial measurements \cite{noh2019tractorbeam}. 

\subsection{MAVLink} MAVLink is the communication protocol used to exchange telemetry, commands, and parameters between the flight controller and external components. In the experimental setup, MAVLink is used to monitor the simulated UAV, read GPS and navigation information, modify simulation parameters, and collect data during the experiments. The attack module is therefore external to the flight controller and interacts with the simulated vehicle through the available MAVLink interfaces. 

\subsection{Waypoint Navigation} A waypoint mission defines a sequence of geographical positions that the UAV must reach. A simplified mission can be represented as \begin{equation} 
W = \{W_0, W_1, \ldots, W_n\} 
\end{equation} 
where each waypoint $W_i$ contains at least a latitude, longitude, and altitude. The flight controller continuously estimates the current vehicle position and generates control actions to move the UAV towards the active waypoint. The mission itself therefore remains unchanged during the experiments presented in this work. Any trajectory deviation is produced by modifying the position information available to the navigation system. 

% ============================================================
% GPS SPOOFING 
% ============================================================ 

\subsection{GPS Spoofing} GPS spoofing is an attack in which an adversary generates or manipulates navigation signals so that a GPS receiver estimates an incorrect position, velocity, or time. Unlike GPS jamming, whose primary objective is to deny access to the navigation signal, spoofing attempts to preserve apparently valid navigation information while modifying its content. The literature commonly distinguishes between hard and soft GPS spoofing \cite{noh2019tractorbeam}. 

\subsubsection{Hard GPS Spoofing} In hard GPS spoofing, the manipulated signal is sufficiently different from the authentic signal that the receiver may lose its original lock before acquiring the spoofed signal. The Tractor Beam paper describes this approach as initially behaving similarly to GPS interference: the receiver may lose its lock and later acquire the spoofed signal. Such a procedure can take a significant amount of time and may activate the target's GPS failsafe mechanism \cite{noh2019tractorbeam}. Hard spoofing is therefore generally unsuitable for experiments in which the objective is to maintain continuous autonomous navigation. 

\subsubsection{Soft GPS Spoofing} 
Soft GPS spoofing, instead, attempts to maintain consistency with the authentic measurements while progressively changing the apparent position. The attacker initially generates measurements close to the authentic position and gradually moves the spoofed position away from the real position. This concept is particularly relevant to the present work because the objective is to investigate whether progressive position manipulation can influence the trajectory without immediately triggering the navigation system's glitch protection. Noh et al. describe this approach as a central component of their Strategy C, where the spoofed position is gradually deviated from the real position while attempting to preserve GPS lock \cite{noh2019tractorbeam}. 

% ============================================================ 
% GPS GLITCH DETECTION 
% ============================================================ 

\subsection{GPS Glitch Detection and Failsafe} GPS measurements cannot always be assumed to be correct. Multipath propagation, signal degradation, receiver errors, and other phenomena may produce significant position errors. ArduCopter therefore implements mechanisms for detecting anomalous GPS measurements. A new GPS position can be compared with a position predicted from the previous position and velocity. The measurement is considered acceptable only when the displacement is compatible with predefined limits. One of the parameters relevant to this mechanism is the maximum allowed GPS glitch radius. Another constraint is related to the maximum physically reasonable position change between consecutive measurements. The EKF provides an additional layer of protection. GPS measurements are compared with the state predicted from inertial sensors. When the difference becomes sufficiently large, the GPS measurement may be rejected and an EKF inconsistency can be reported. This behaviour is fundamental for the experiments in this work because increasing the GPS offset too rapidly does not necessarily result in a corresponding change of the UAV trajectory. Instead, the flight controller may classify the measurement as anomalous and reject it. Therefore, the effective manipulation is not determined only by the absolute GPS offset, but also by its temporal evolution. 

% ============================================================ 
% SAFE HIJACKING 
% ============================================================ 

\subsection{Safe Hijacking} The concept of safe-hijacking was introduced by Noh et al. in the context of autonomous consumer drones. Instead of simply disrupting a UAV, safe-hijacking aims to influence its autonomous behaviour so that the vehicle can be moved away from a sensitive area \cite{noh2019tractorbeam}. The authors analyzed several consumer UAVs and classified them according to their GPS failsafe behaviour. They then designed different spoofing strategies depending on the characteristics of each class. The key observation is that a GPS manipulation cannot be designed independently from the target's navigation and failsafe mechanisms. A spoofed position that is too different from the expected position may trigger a failsafe, preventing the manipulation from producing the desired trajectory. The paper therefore distinguishes multiple strategies. In particular, Strategy C targets systems for which entering a GPS-independent failsafe mode prevents subsequent GPS manipulation from influencing the vehicle. For these systems, the spoofed GPS position must remain sufficiently consistent with the authentic measurements to avoid triggering the failsafe mechanism \cite{noh2019tractorbeam}. 

\subsubsection{Relation to ArduCopter} An important aspect of the Tractor Beam study is its analysis of the 3DR Solo, whose software was based on ArduCopter. The authors analyzed both the EKF failure detection mechanism and the ArduCopter path-following algorithm. Their analysis showed that ArduCopter uses an intermediate target position during waypoint navigation. The Intermediate Target Position (ITP) is advanced along the trajectory toward the next waypoint, and the vehicle attempts to move towards this target position. The paper further describes a ``leash'' mechanism that limits how far the vehicle can deviate from the planned track before the intermediate target is updated \cite{noh2019tractorbeam}. This behaviour is particularly relevant to the present study. If the flight controller believes that the UAV is located at a position different from its physical position, the path-following algorithm may generate control actions intended to return the vehicle to the planned trajectory. Consequently, GPS manipulation can potentially influence the physical trajectory indirectly, without directly issuing movement commands to the flight controller. 

\subsubsection{Adaptive GPS Manipulation} The strategy investigated in this work follows the same general principle of adaptive manipulation. Instead of applying a single large GPS offset, the spoofed position is progressively modified. At each iteration, the current GPS state is observed and the next manipulation is calculated according to: 

\begin{equation} P_{\mathrm{spoof}}(t) = P_{\mathrm{real}}(t) + \Delta P(t) \end{equation} 
where 
\begin{equation} \Delta P(t) = \begin{bmatrix} \Delta \phi(t)\\ \Delta \lambda(t) \end{bmatrix} \end{equation} 

represents the accumulated GPS manipulation. The objective is to keep the change between consecutive measurements small enough to avoid immediately activating the GPS glitch protection, while progressively changing the position perceived by the navigation system. The general control loop can therefore be represented as: 

\begin{equation} \text{Real position} \rightarrow \text{GPS manipulation} \rightarrow \text{Estimated position} \rightarrow \text{Flight controller} \rightarrow \text{Physical trajectory} \end{equation} 

The resulting physical trajectory changes the subsequent real position, which is then used to calculate the next spoofed position. This creates a feedback loop in which the manipulation is continuously adapted to the current state of the UAV.